Al analizar la evolución del precio de cierre utilizando un índice base 100 se observa que las empresas con mayor rendimiento acumulado siguieron trayectorias muy diferentes. BNTX presentó el comportamiento más volátil del grupo, alcanzando un máximo cercano a 3,140 puntos durante agosto de 2021 para posteriormente experimentar una fuerte corrección que redujo su índice hasta aproximadamente 718 puntos al finalizar el período analizado. Este comportamiento refleja un crecimiento acelerado seguido de una importante pérdida de valor relativo respecto a su máximo histórico.
Por su parte, TSLA también registró una elevada volatilidad. Su índice alcanzó un máximo cercano a 1,803 puntos durante noviembre de 2021; sin embargo, posteriormente presentó una corrección importante que redujo su desempeño hasta finalizar el período alrededor de 1,090 puntos. A pesar de esta disminución, la empresa mantuvo uno de los mayores rendimientos acumulados del conjunto analizado, evidenciando una apreciación considerable del precio de sus acciones durante el horizonte de estudio.
En contraste, NVIDIA mostró una trayectoria de crecimiento más estable y sostenida. Aunque registró fluctuaciones propias del mercado, su comportamiento fue menos pronunciado que el observado en BNTX y TSLA. Además, finalizó el período con un índice cercano a 1,232 puntos y un máximo de aproximadamente 1,291 puntos, lo que sugiere una apreciación más consistente del precio de la acción y una menor dependencia de movimientos especulativos de corto plazo.